% Part: first-order-logic
% Chapter: syntax-and-semantics
% Section: covered-structures

\documentclass[../../../include/open-logic-section]{subfiles}

\begin{document}

\olfileid[fr]{fol}{syn}{cov}

\olsection{\printtoken{P}{structure} couvertes pour les langages du premier ordre}

\begin{explain}
Rappelons qu'un terme est \emph{clos} s'il ne contient aucune
!!{variable}.
\end{explain}

\begin{defn}[Valeur d'un terme clos]
Si $t$ est un terme clos du langage~$\Lang L$ et $\Struct M$ une
!!{structure} pour~$\Lang L$, la \emph{!!{value}}~$\Value{t}{M}$
est définie comme suit :
\begin{enumerate}
\item Si $t$ est simplement la !!{constant}~$c$, alors
  $\Value{c}{M} = \Assign{c}{M}$.
\item Si $t$ est de la forme $\Atom{f}{t_1, \ldots, t_n}$, alors
  \[
  \Value{t}{M} = \Assign{f}{M}(\Value{t_1}{M}, \ldots,
  \Value{t_n}{M}).
  \]
\end{enumerate}
\end{defn}

\begin{defn}[Structure couverte]
Une !!{structure} est \emph{couverte} si tout !!{element} de son
!!{domain} est la !!{value} d'un terme clos.
\end{defn}

\begin{ex}
Soit ~$\Lang L$ le langage qui contient les !!{constant}s
$\Obj{zero}$, $\Obj{one}$, $\Obj{two}$, \dots, le !!{predicate}
binaire~$<$ et les !!{function}s binaires $+$ et~$\times$.
Considérons la !!{structure}~$\Struct M$ pour~$\Lang L$ dont le
domaine est $\Domain M = \{0, 1, 2, \ldots \}$ et où
$\Assign{\Obj{zero}}{M} = 0$,
$\Assign{\Obj{one}}{M} = 1$, $\Assign{\Obj{two}}{M} = 2$, et ainsi
de suite. Le symbole de relation binaire~$<$ est interprété par
l'ensemble $\Assign{<}{M}$ de tous les couples
$\tuple{c_1, c_2} \in \Domain{M}^2$ tels que $c_1$ est inférieur
à~$c_2$ : par exemple,
$\tuple{1, 3} \in \Assign{<}{M}$, mais
$\tuple{2, 2} \notin \Assign{<}{M}$. Pour le !!{function} binaire~$+$,
on définit $\Assign{+}{M}$ de la manière usuelle : ainsi,
$\Assign{+}{M}(2,3)$ vaut~$5$ ; il en va de même pour le
!!{function} binaire~$\times$. La !!{value} de $\Obj{four}$ est
donc~$4$, et celle de
$\times(\Obj{two}, +(\Obj{three},\Obj{zero}))$, que l'on peut aussi
écrire sous forme infixe
$\Obj{two} \times (\Obj{three} + \Obj{zero})$, est
\begin{multline*}
\Value{\times(\Obj{two}, +(\Obj{three},\Obj{zero}))}{M} =\\
\begin{aligned}
& \Assign{\times}{M}(\Value{\Obj{two}}{M}, \Value{+(\Obj{three}, \Obj{zero})}{M})\\
& = \Assign{\times}{M}(\Value{\Obj{two}}{M}, \Assign{+}{M}(\Value{\Obj{three}}{M},
\Value{\Obj{zero}}{M})) \\
& = \Assign{\times}{M}(\Assign{\Obj{two}}{M}, \Assign{+}{M}(\Assign{\Obj{three}}{M},
\Assign{\Obj{zero}}{M})) \\
& = \Assign{\times}{M}(2, \Assign{+}{M}(3, 0)) \\
& = \Assign{\times}{M}(2, 3) \\
& = 6
\end{aligned}
\end{multline*}
\end{ex}

\begin{prob}
La !!{structure}~$\Struct N$, modèle standard de l'arithmétique,
est-elle couverte ? Expliquer.
\end{prob}

\end{document}
